The feast is ready. That is where today's Gospel begins. Before the guests
have answered, before anyone has taken a place, the king has prepared the
wedding feast for his son. The invitation asks them to receive a gift.
They have not bought the food or earned the celebration. They are being
called to share the king's joy.

Yet some have other things to do. A farm needs attention. Business calls.
Others go further and attack the messengers. The king judges the murderers,
but their refusal does not defeat his purpose. He sends his servants out
again, into the roads, to gather everyone they find. At last the hall is
full. Good and bad have come together.

Then the king enters, and our attention falls on one guest. He is already
inside. He has answered the invitation. But he has no wedding garment,
and he is cast out. The ending is severe. Jesus brings us from the joy
of being invited to the seriousness of how we receive the invitation.
Being in the room does not finish the question.

What, then, is this garment? Saint Augustine and Saint Gregory the Great
answer: charity. The love of God and the love of our neighbour. Gregory
explains why that garment belongs to this wedding. Christ came in love
to unite the Church to himself. To celebrate his union while refusing
his love is to refuse the very good we have come to celebrate.

The question reaches our own hearts. We can be present among Christ's
people, know the words of faith, and value the Church's worship, while
keeping some part of our life closed to charity. Augustine asks whether
possessing great religious gifts is enough. His answer is searching:
the gifts are good, but we must let them bear fruit in us. Coming to
the feast calls us to love the Bridegroom and the people he gathers.

This love reaches the ordinary places where we spend our days. Today's
Epistle begins with the renewal of the mind and ends with hands that
have something to give. Between those two points, Saint Paul speaks
about truth, anger and work. He gives charity a voice and a pair of hands.

First, truth. We are members of one another. Someone else has to live
with the consequences of what we say. Saint John Chrysostom makes this
plain through the image of a body. If the eye sees danger and hides it
from the foot, the whole body suffers. In the same way, a lie damages
the common life to which the liar belongs. Truthfulness serves the
person who must depend on our word.

Then anger. Paul places the setting sun beside our resentment. Chrysostom
notices how anger can grow in the quiet of the night. An injury is
rehearsed, an accusation repeated inwardly, until everything the other
person does seems to confirm the charge. The heart becomes more willing
to believe a slander because it already wants the neighbour condemned.

There is something here we can begin to do today. A wrong may still
need correcting. Its consequences may take time to repair. But we can
stop giving resentment another meal. We can refuse to repeat the
accusation for the pleasure of repeating it. We can seek peace without
pretending that the injury never happened. Charity begins to change
what we do with the injury.

Finally, work. Paul tells the former thief to stop stealing and to
labour at what is good, so that someone in need can receive help.
The person once harmed a neighbour by taking; now he can help a
neighbour by giving. The appointed reading ends with that other
person's need. Our work and our possessions have a purpose beyond
ourselves. Charity asks whose need can be met by what we have.

This does not measure anyone's worth by strength or earnings. The
person who needs help is the person Paul tells us to serve. Poverty
and weakness do not cancel the invitation. The question concerns
what we can give and what we choose to withhold. Even the gift
itself asks something of the heart: do we seek the neighbour's good,
or chiefly the satisfaction of being admired?

Now the Gospel's garment is very close to us. It touches the next
truthful answer, the resentment we stop nursing, the help we can
actually give. And it reaches the neighbour who is hardest to love.
Augustine asks us to desire an enemy's healing. We can oppose the
evil someone does and still desire that the person be freed from it.
Love seeks his conversion rather than taking delight in his ruin.

That is a demanding love. Its pattern is Christ himself, the
Bridegroom who gave himself for us. Augustine directs our attention
to Christ upon the cross, praying for his enemies. The garment's
meaning becomes clearest there. The Lord does not ask us for a love
that he has withheld from us. He draws us into his own self-giving.

Where can such love grow in people as weak as we are? Listen to
the prayers of this Mass. The Collect asks God to remove what
obstructs us, so that body and mind may serve him freely. We ask
for the freedom to answer him. Our conversion begins within his
mercy, and every real step forward still needs his help.

The prayer over the gifts asks that the offering bring saving
fruit to those who share in it. We bring our need before God
even as we offer. Christ's gift is real; our hearts still need
to receive it fruitfully. The Gospel's question follows us to
the altar: will the love we receive become the life we live?

The Communion chant places God's command beside a longing for
our own ways to be directed toward obedience. Augustine hears
that longing as prayer. Knowing what is right does not give us
all the strength to do it. We ask for help because we want to
walk the way God shows us. Dependence on grace can become the
beginning of action: the truthful word spoken, the revenge
refused, the gift placed in another's hands.

After Communion, the Church asks for healing. God's medicine
is to free us from our perversity and keep us faithful to his
commands. The prayer reaches beneath an isolated act to the
desire from which it comes. We need to be healed where we
prefer the lie, enjoy the grievance, or hold back what another
person needs. The Eucharist brings us to Christ, whose grace
can change that inward life and sustain it.

Gregory gives a word of hope to the person who hears this and
sees how much remains unfinished. To have charity imperfectly
is different from having none at all. Weak love needs growth.
The struggle to love more faithfully is not a reason to despair
of God's mercy. The question is whether we will receive his
healing and let our love grow.

So take one concrete act of charity into the day ahead. Tell
the truth where someone is relying on you. Interrupt the
grievance you have been feeding. Give help that meets a real
need. Choose the act that answers the place where you are
resisting love. Ask for God's help, and begin. Such a beginning
does not buy the invitation. It answers the love that has
already called you.

The king's judgment remains serious. We cannot postpone our
answer for ever. But the feast remains the feast of his Son,
and its promise is greater than our present weakness. Christ
calls his people toward the joy in which charity is perfected,
where love no longer struggles against the evil that divides us.
The feast is ready. Receive the invitation with gratitude.
Let the Bridegroom's love take hold of your heart, and let
that love become a life.
